% !TeX program = xelatex
\documentclass[11pt,a4paper]{article}
\usepackage{fontspec}
\newfontfamily\greekfont[Script=Greek]{Walleye}
\newfontfamily\greekfontsf[Script=Greek]{Walleye}
\usepackage{polyglossia}
\setdefaultlanguage{greek}
\usepackage{enumitem}
\usepackage{amssymb,amsmath,amsthm,changepage,extarrows,subcaption,tikz,xcolor,xfrac}
\usepackage[margin=1in]{geometry}
\usepackage{mathpazo}
\usepackage[colorlinks=true, urlcolor=blue]{hyperref}

\usetikzlibrary{calc}
\usetikzlibrary{intersections}
\usetikzlibrary{patterns}
\usetikzlibrary{decorations.footprints}

\pgfdeclarelayer{bg}
\pgfsetlayers{bg,main}

\renewcommand{\thesubfigure}{\roman{subfigure}}

\newcommand{\abs}[1]{{\left|#1\right|}}
\newcommand{\pn}[1]{{\left(#1\right)}}
\newcommand{\qn}[1]{{\left[#1\right]}}
\newcommand{\mn}[1]{{\left\{#1\right\}}}
\newcommand{\R}{\mathbb{R}}
\newcommand{\Q}{\mathbb{Q}}
\newcommand{\Z}{\mathbb{Z}}
\newcommand{\N}{\mathbb{N}}
\newcommand{\tor}{{\text{ ή }}}
\newcommand{\tand}{{\text{ και }}}
\newcommand{\tif}{{\text{ αν }}}

\definecolor{twitter}{RGB}{29,161,242}
\definecolor{wordpress}{RGB}{71,167,200}
\definecolor{instagram}{RGB}{240,76,92}
\definecolor{facebook}{RGB}{66,103,178}

\let\sin\relax
\DeclareMathOperator{\sin}{\text{{ημ}}}

\let\cos\relax
\DeclareMathOperator{\cos}{\text{{συν}}}

\let\tan\relax
\DeclareMathOperator{\tan}{\text{{εφ}}}

\let\cot\relax
\DeclareMathOperator{\cot}{\text{{σφ}}}

\theoremstyle{definition}

\newtheorem{them}{Θέμα}
\newtheorem{prop}{Πρόταση}[section]
\newtheorem{cor}{Πόρισμα}[section]

\everymath{\displaystyle}

\title{\Huge{\bf Λάμπες\ldots}\\{\Large'Η πώς να σκοτώσετε την ώρα σας στην καραντίνα}}
%\author{{\color{facebook}\faFacebook}\ \href{www.facebook.com/aftermaths.gr}{\en{/aftermaths.gr}} {\color{instagram}\faInstagram}\ \href{https://www.instagram.com/aftermathsgr}{\en{@aftermathsgr}} {\color{twitter}\faTwitter}\,\href{https://twitter.com/aftermathsgr}{\en{@aftermathsgr}}}
\date{\today}

\begin{document}
	\maketitle
	\section{Ευκολοφώτιστα δωμάτια}
	Ας πούμε ότι έχουμε ένα δωμάτιο σαν αυτό που φαίνεται στο Σχήμα \ref{fig:1}.
	
	\begin{figure}[!htb]
		\centering
		\begin{tikzpicture}
		\draw[thick] (0,0) rectangle (6,4);
		\end{tikzpicture}
		\caption{'Ενα βαρετό δωμάτιο.}
		\label{fig:1}
	\end{figure}

	Εντάξει, δεν έχει σχεδόν τίποτα μέσα, αλλά δε μας νοιάζει αυτό, προς το παρόν. Θέτουμε το εξής ερώτημα: Πού πρέπει να τοποθετήσουμε μία λάμπα έτσι ώστε να φωτίζεται όλο το δωμάτιο;
	
	Σιγά τ' αυγά. 'Οπου κι αν θέλει η ψυχή μας να βάλουμε τη λάμπα, αν αυτή είναι αρκετά δυνατή κι όχι κάποιο μικρό κι αδύναμο λαμπάκι, θα μπορεί εύκολα να φωτίσει όλο το δωμάτιο. Μάλιστα,το ίδιο ισχύει και για το πιο φουτουριστικό δωμάτιο που φαίνεται στο Σχήμα \ref{fig:2}.
	
	\begin{figure}[!htb]
		\centering
		\begin{tikzpicture}
		\draw[thick] (0,0) -- (2,-1) arc (0:atan(-1.5):3cm) arc (-90:-180:2cm) -- cycle;
		\end{tikzpicture}
		\caption{'Ενα βαρετά φουτουριστικό δωμάτιο.}
		\label{fig:2}
	\end{figure}

	Αμφότερα τα δωμάτια που φαίνονται στα Σχήματα \ref{fig:1} και \ref{fig:2} έχουν ένα κοινό χαρακτηριστικό: \textit{είναι κυρτά}. Στην πράξη, ένα σύνολο είναι κυρτό αν μπορούμε κάθε δύο σημεία του να τα ενώσουμε με ένα ευθύγραμμο τμήμα έτσι ώστε όλο το ευθύγραμμο τμήμα να παραμένει \textit{εντός} του συνόλου. Πιο απλά, αν στεκόμενοι σε ένα οποιοδήποτε σημείο του δωματίου μπορούμε να δούμε κάθε άλλο σημείο του, τότε αυτό είναι κυρτό. 'Ετσι, για παράδειγμα, ενώ τα δωμάτια στα Σχήματα \ref{fig:1} και \ref{fig:2} είναι κυρτά, το δωμάτιο που φαίνεται στο Σχήμα \ref{fig:3} δεν είναι, καθώς αν τοποθετήσουμε την Αλίκη στη θέση $A$ και τον Βασίλη\footnote{Κατ' αναλογία του διάσημου ζευγαριού Alice \& Bob που στοιχειώνει εμμονικά τόσα και τόσα παραδείγματα σε μυριάδες επιστημονικά και μη συγγράμματα, ειδικά στον χώρο της πληροφορικής.} στη θέση $B$ τότε δε θα μπορούν να ανταλλάξουν ούτε ένα βλέμμα!
	
	\begin{figure}[!htb]
		\centering
		\begin{tikzpicture}
		\draw[thick] (0,0) -- (4,0) -- (4,-4) -- (0,-4) -- (2,-2) -- cycle;
		\node[circle, inner sep=1.5pt, draw=black, fill=black, label={right}:{$A$}](a) at (1.5,-.75){};
		\node[circle, inner sep=1.5pt, draw=black, fill=black, label={right}:{$B$}](b) at (1.5,-3.25){};
		\end{tikzpicture}
		\caption{Χωρίζοντας την Αλίκη και τον Βασίλη.}
		\label{fig:3}
	\end{figure}

	Αυστηρά μιλώντας, αν $S\subseteq\R^2$ είναι ένα σχήμα στο επίπεδο τότε θα λέμε ότι αυτό είναι \textit{κυρτό} αν για κάθε $x,y\in S$ και κάθε $\lambda\in(0,1)$ ισχύει και ότι:
	\begin{equation}\label{eq:convex}
	\lambda x+(1-\lambda)y\in S
	\end{equation}
	Αυτός είναι ο άχαρος τρόπος να περιγράψουμε στα μαθηματικένια το «Η Αλίκη κι ο Βασίλης μπορούν να βλέπονται όπου κι αν σταθούν στο $S$». Αξίζει, ωστόσο, να σταθούμε λίγο στη διαίσθηση πίσω από αυτόν τον ορισμό. Μπορούμε να ξαναγράψουμε την Εξίσωση \eqref{eq:convex} ως εξής:
	\begin{equation}\label{eq:convex_alt}
	y+\lambda(x-y)\in S.
	\end{equation}
	Πρακτικά, απλώς κάναμε κάτι επιμεριστικές και συμμαζέψαμε μαζί τα $\lambda$. Ωστόσο, από την Εξίσωση \eqref{eq:convex_alt} μπορούμε πιο εύκολα να «διαβάσουμε» πώς ο ορισμός που γράψαμε παραπάνω πράγματι αποτυπώνει τη διαισθητική ιδέα μας για τα διασταυρούμενα βλέμματα της Αλίκης και του Βασίλη. Ας δούμε λίγο το Σχήμα \ref{fig:4} στο οποίο έχουμε σχεδιάσει δύο σημεία $x,y$ και διάφορα βελάκια.
	
	\begin{figure}[!htb]
		\centering
		\begin{tikzpicture}
			\draw[thick, ->] (-1,0) -- (6,0);
			\draw[thick, ->] (0,-1) -- (0,4);
			\node[circle, inner sep=1.5pt, draw=black, fill=black, label={left}:{$y$}](y) at (1,2){};
			\node[circle, inner sep=1.5pt, draw=black, fill=black, label={right}:{$x$}](x) at (4.5,3.5){};
			\node[circle, inner sep=1.5pt, draw=black, fill=black](z) at (3.8,3.2){};
			\draw[thick, red!50, ->] (y) -- (z)node[pos=.5, yshift=8pt, rotate=atan(3/7)]{$\lambda(x-y)$};
			\draw[thick, red!50, ->] (0,0) -- (y);
			\draw[thick, blue!50, ->] (0,0) -- (z)node[pos=.55, yshift=-10pt, rotate=atan(7/9)]{$y+\lambda(x-y)$};
			\draw[dashed, ->] (y) -- (x);
		\end{tikzpicture}
		\caption{Εξηγώντας τον ορισμό της κυρτότητας.}
		\label{fig:4}
	\end{figure}

	'Οπως ίσως θα θυμόμαστε από το σχολείο, τα σημεία του επιπέδου μπορούμε να τα σκεφτόμαστε είτε ως\ldots σημεία (κουκκίδες) είτε ως διανύσματα (βελάκια) με αρχή την αρχή των αξόνων (δηλαδή, το σημείο $O(0,0)$). Τώρα, ανακαλούμε επίσης από τα σχολικά μας χρόνια και το γεγονός ότι η διαφορά δύο διανυσμάτων είναι επίσης διάνυσμα με αρχή το πέρας του «αφεραιτέου» και πέρας το πέρας «μειωτέου». Στο σχήμα μας, αυτό σημαίνει ότι το $x-y$ αντιστοιχεί στο διακεκομμένο διάνυσμα που ξεκινάει από το σημείο $y$ και καταλήγει στο σημείο $x$. Τώρα, το $\lambda(x-y)$ είναι ένα διάνυσμα που η μόνη διαφορά του από το $x-y$ είναι το μήκος του, το οποίο είναι το μήκος του $x-y$ πολλαπλασιασμένο επί $\lambda$. Επειδή, στην περίπτωσή μας, $\lambda\in(0,1)$, το μήκος του $\lambda(x-y)$ είναι πάντοτε μικρότερο αυτού του $x-y$, και άρα πάντοτε έχουμε μία εικόνα όπως αυτή που αποτυπώνεται στο Σχήμα \ref{fig:4}.
	
	Τώρα, αν θυμηθούμε ότι το άθροισμα δύο διανυσμάτων είναι απλώς ένα διάνυσμα με αρχή την αρχή του πρώτου και πέρας το πέρας του δεύτερου (υπό την προϋπόθεση το πέρας του πρώτου να είναι η αρχή του δεύτερου), τότε εύκολα μπορούμε να ερμηνεύσουμε αυτό που βλέπουμε στο Σχήμα \ref{fig:4}. Το μπλε διάνυσμα είναι ακριβώς το ίδιο με το άθροισμα των δύο κόκκινων διανυσμάτων ή, με άλλα λόγια, οι Εξισώσεις \eqref{eq:convex} και \eqref{eq:convex_alt} μας λένε το ίδιο πράγμα με άλλα λόγια. Επομένως, πράγματι, το σημείο $y+\lambda(x-y)$ είναι πάντοτε ένα σημείο επί του ευθυγράμμου τμήματος με άκρα τα $x,y$ (για $0<\lambda<1$) και άρα πράγματι ο μαθματικένιος ορισμός της κυρτότητας μάς λέει αυτό που μας λένε κι η Αλίκη με τον Βασίλη με τα μάτια τους.
	
	\section{Φωτίζοντας τις σκιές\ldots}
	Είπαμε ότι αν έχουμε ένα κυρτό δωμάτιο μπορούμε πολύ εύκολα με μία αρκετά δυνατή λάμπα να το φωτίσουμε όλο τοποθετώντας τη λάμπα όπου θέλουμε. Αν όμως το δωμάτιό μας δεν είναι κυρτό αλλά είναι, για παράδειγμα το δωμάτιο του Σχήματος \ref{fig:3}; Εκεί υπάρχουν θέσεις όπου μία λάμπα δεν μπορεί να φωτίσει όλο το δωμάτιο (όπως η θέσεις $A$ και $B$) και θέσεις όπου η λάμπα μπορεί να το φωτίσει ολόκληρο (μπορείτε να βρείτε μία τέτοια;). Μπορούμε, μάλιστα, να πάρουμε ένα ακόμα πιο «άσχημο» δωμάτιο, όπως αυτό που φαίνεται στο Σχήμα \ref{fig:5}.
	
	\begin{figure}[!htb]
		\centering
		\begin{tikzpicture}
		\draw[thick] (0,0) -- (0,4) -- (4,4) -- (4,0) -- (3,0) -- (3,2.5) -- (1,2.5) -- (1,0) -- cycle;
		\end{tikzpicture}
		\caption{'Ε´να πραγματικά παράξενο δωμάτιο.}
		\label{fig:5}
	\end{figure}
	
	Σε αυτό το δωμάτιο δεν υπάρχει θέση στην οποία με μία λάμπα να μπορούμε να το φωτίσουμε. Σαφώς, κάποιος μπορεί να πει ότι, εντάξει, αν παίξουμε λίγο με το χρώμα των τοίχων αλλά και με κάποιον καθρέφτη, όπου και να βάλουμε τη λάμπα κάπως λύνεται το πρόβλημα. Αλλά, δεν είναι αυτή η πρόθεσή μας. Αυτό θα ήταν μία πολύ πρακτική αντιμετώπιση του ζητήματος κι εδώ έχουμε αποφασίσει να κάνουμε μαθηματικά. Κάνουμε, λοιπόν, τις εξής υποθέσεις:
	\begin{itemize}
		\item έχουμε στη διάθεσή μας μονάχα μία λάμπα.
		\item Η λάμπα μας είναι όσο ισχυρή θέλουμε.
		\item Το φως της λάμπας δεν αντανακλάται (επαρκώς) στους τοίχους έτσι ώστε να πούμε ότι κάποιος χώρος φωτίζεται (επαρκώς) από ανακλώμενες ακτίνες φωτός στους τοίχους.
	\end{itemize}
	'Ετσι, για παράδειγμα, αν τοποθετήσουμε μία λάμπα στη θέση $A$, όπως φαίνεται στο Σχήμα \ref{fig:6}, τότε αυτή θα φωτίζει επαρκώς αποκλειστικά και μόνο τη χρωματισμένη περιοχή.
	\begin{figure}[!htb]
		\centering
		\begin{tikzpicture}
		\fill[red!10] (0,13/6) -- (1,2.5) -- (3,2.5) -- (4,1.5) -- (4,4) -- (0,4) -- cycle;
		\draw[thick] (0,0) -- (0,4) -- (4,4) -- (4,0) -- (3,0) -- (3,2.5) -- (1,2.5) -- (1,0) -- cycle;
		\node[circle, inner sep=1.5pt, draw=black, fill=black, label={right}:{$A$}](a) at (2.5,3){};
		\draw[dashed] (a) -- (4,1.5);
		\draw[dashed] (a) -- (0,13/6);
		\begin{scope}
		\clip (0,13/6) -- (1,2.5) -- (3,2.5) -- (4,1.5) -- (4,4) -- (0,4) -- cycle;
		\foreach \i in {0,10,...,360} {
			\draw[thin, blue!30] (a) -- ($(a)+(\i:3)$);
		}
		\end{scope}
		\end{tikzpicture}
		\caption{Φέξε μου και γλίστρησα.}
		\label{fig:6}
	\end{figure}
	Επομένως, ένα εύλογο ερώτημα γεννιέται: Πώς μπορούμε να βρούμε μία θέση από την οποία η λάμπα θα φωτίζει το μεγαλύτερο δυνατό μέρος του δωματίου;
	
	Αρχικά, για να απαντήσουμε το παραπάνω ερώτημα καλό θα ήταν να προσδιορίσουμε για τι είδους δωμάτια μιλάμε. Αν βασιστούμε στην κοινή μας εμπειρία περί δωματίων, τότε ίσως θα συμφωνήσουμε ότι ένα δωμάτιο μπορεί να περιγραφεί αρκετά καλά από ένα  πολύγωνο. Εντάξει, σίγουρα υπάρχουν δωμάτια που οι τοίχοι τους δεν είναι όλοι επίπεδοι (ίσως όπως στο Σχήμα \ref{fig:2}) αλλά, πόσα τέτοια δωμάτια να υπάρχουν; Ας ελπίσουμε, για καλή μας τύχη, όχι και πολλά\ldots
	
	Τώρα, το επόμενο βήμα που πρέπει να κάνουμε είναι να περιγράψουμε αυστηρά πώς προκύπτει η ροζ περιοχή στο Σχήμα \ref{fig:6}. Με άλλα λόγια, αν μας δωθούν ένα πολύγωνο, $P$, κι ένα σημείο, $A$, στο εσωτερικό του, πώς μπορούμε να προσδιορίσουμε ποια θα είναι η περιοχή που φωτίζει μία λάμπα, αν τοποθετηθεί εκεί; Μία απλή σκέψη είναι να σκεφτούμε τη λάμπα σαν ένα μεγάλο «μπεκ» αυτόματου ποτίσματος που «ποτίζει» το δωμάτιο με φως (τύφλα νά' χει ο Παλαμάς). 'Ετσι, για να προσδιορίσουμε τη φωτιζόμενη περιοχή, μπορούμε να αρχίσουμε να τραβάμε, με αρχή το σημείο $A$ ευθύγραμμα τμήματα μέχρι να πέσουμε πάνω σε κάποια κορυφή ή πλευρά του $P$. Η περιοχή που θα «ζωγραφιστεί» από όλα αυτά τα ευθύγραμμα τμήματα είναι και αυτή που μας ενδιαφέρει.
	
	'Εχοντας στον νου τα παραπάνω, ένας εύλογος τρόπος να ορίσουμε το υποσύνολο του $P$ που φωτίζεται από το $A$ είναι ο εξής:
	\begin{equation}\label{eq:light_pa}
	\phi_P(A):=\{x\in P:[A,x]\subseteq P\},
	\end{equation}
	όπου ως $[A,x]$ ορίζουμε το ευθύγραμμο τμήμα με άκρα τα σημεία $A$ και $x$, δηλαδή:
	\begin{equation}\label{eq:segment}
	[A,x]:=\{\lambda A+(1-\lambda)x:\lambda\in[0,1]\}.
	\end{equation}
	Να πούμε κάπου εδώ ότι θεωρούμε ότι οι κορυφές και οι πλευρές ενός πολυγώνου είναι μέρος του πολυγώνου, οπότε στην προκειμένη, οι διακεκομμένες γραμμές βρίσκονται εξ ολοκλήρου στη φωτιζόμενη περιοχή, $\phi_P(A)$. Εδώ μπορεί κανείς να εγείρει διάφορες ενστάσεις ως προς το αν το φως πράγματι περνά ή όχι από τις γωνίες κι αν θα έπρεπε, όπως φαίνεται στο Σχήμα \ref{fig:7} να συμπεριλάβουμε τις διακεκομμένες γραμμές στο $\phi_P(A)$ ή όχι. Η απάντηση που λακωνικά δίνουμε είναι η εξής: «Δε μας νοιάζει». Οι κορυφές του $P$ είναι πεπερασμένες στο πλήθος, επομένως και τα ευθύγραμμα τμήματα που θα μας απασχολήσουν κατ' αυτόν τον τρόπο είναι επίσης πεπερασμένα στο πλήθος. 'Αρα, από άποψη εμβαδού, δε συνεισφέρουν κάπως στην $\phi_P(A)$. Συνακόλουθα, είτε τα συνυπολογίσουμε είτε όχι, δεν έχει καμία, μα καμία, διαφορά - και άρα, είναι «νόμιμο» να τα λάβουμε κι αυτά υπόψιν για να μην μπλέκουμε ή να τα αγνοούμε όποτε δε μας βολεύει να τα βλέπουμε.
	
	\begin{figure}[!htb]
		\centering
		\begin{tikzpicture}
		\fill[red!10] (0,13/6) -- (1,2.5) -- (3,2.5) -- (4,1.5) -- (4,4) -- (0,4) -- cycle;
		\draw[thick] (0,0) -- (0,4) -- (4,4) -- (4,0) -- (3,0) -- (3,2.5) -- (1,2.5) -- (1,0) -- cycle;
		\node[circle, inner sep=1.5pt, draw=black, fill=black, label={right}:{$A$}](a) at (2.5,3){};
		\draw[dashed] (3,2.5) -- (4,1.5);
		\draw[dashed] (1,2.5) -- (0,13/6);
		\end{tikzpicture}
		\caption{Το μετράει; Δεν το μετράει; Το μετράει; Δεν το μετράει; \ldots}
		\label{fig:7}
	\end{figure}

	Με αυτόν τον ορισμό κατά νου, αυτό που αναζητούμε για δεδομένο πολύγωνο $P$ είναι ένα σημείο $A$ έτσι ώστε το ακόλουθο ολοκλήρωμα να παίρνει τη μέγιστη τιμή του:
	\begin{equation}\label{eq:max_int}
	\iint_{\phi_P(A)}dxdy.
	\end{equation}
	Εντάξει, τι τα θέλαμε τώρα τα ολοκληρώματα; Από εκεί που μιλούσαμε όμορφα κι ωραία για λάμπες και φωτάκια, ξαφνικά, νά' σου ένα διπλό ολοκλήρωμα από το πουθενά. Η αλήθεια είναι ότι τα ολοκληρώματα είναι πολύ βαριά εργαλεία και ίσως το να καταφύγουμε αμέσως σε αυτά να μην είναι και η καλύτερη λύση\footnote{Η, πιο απλά, τώρα δε βλέπουμε κάποιον σαφή τρόπο να βρούμε τη μέγιστη τιμή του ολοκληρώματος αυτού ως προς $A$ και άρα προσπαθούμε να ξεγλιστρήσουμε από τα χέρια του απαιτητικού αναγνώστη\ldots}. Πριν, λοιπόν, καταλήξουμε να κάνουμε βαρύ απειροστικό λογισμό, ας προσπαθήσουμε να μελετήσουμε λίγο διάφορες ιδιότητες του\footnote{'Ισως παρατηρήσατε ότι μία λέμε «η $\phi_P(A)$ και μία «το $\phi_P(A)$». Αυτό το παθαίνουμε συχνά στα μαθηματικά, καθώς τα περισσότερα πράγματα έχουν ένα γραμματικό γένος από το πλαίσιο που ορίζονται (εδώ, «\textit{η φωτιζόμενη περιοχή} $\phi_P(A)$»{)} και τουλάχιστον ένα από το ευρύτερο μαθηματικό πλαίσιο στο οποίο υπάγονται (εδώ «\textit{το σύνολο} $\phi_P(A)$»{)}. Αυτό το gender fluidity χαρακτηρίζει εγγενώς τα μαθηματικά, οπότε θα το ανεχτείτε και στη συνέχεια.} $\phi_P(A)$.
	
	\section{Φωτεινά πολύγωνα κι άλλα παράξενα}
	Μία βασική παρατήρηση που μπορούμε να κάνουμε για το σύνολο $\phi_P(A)$ είναι ότι, εν γένει, δεν είναι κυρτό - κάτι που εύκολα επαληθεύεται κι από το Σχήμα \ref{fig:6}. Ωστόσο, από τον ορισμό του μπορούμε να δούμε ότι είναι ένα πολύγωνο που περιέχεται στο $P$ και περιέχει το $A$ (έχει κι άλλες ιδιότητες, αλλά baby steps, που λέμε και στο χωριό μου). Το γεγονός ότι είναι πολύγωνο ίσως είναι σχετικά προφανές, αλλά αξίζει τον κόπο να προσπαθήσουμε να το «αποδείξουμε» κάπως πιο αυστηρά, κυρίως γιατί τρώγοντας έρχετ' η όρεξη και γιατί, όσο διαπραγματευόμαστε τις έννοιες που μας απασχολούν αποκτάμε και μία μεγαλύτερη οικειότητα μαζί τους (μία οικειότητα που ελπίζουμε να μας βοηθήσει να τις καταλάβουμε καλύτερα).
	
	Ας πάρουμε, λοιπόν, ένα πολύγωνο $P$ κι ένα σημείο $A$ στο $P$. Ο τρόπος που κατασκευάζουμε το $\phi_P(A)$ είναι αρκετά απλός: ξεκινώντας από το $A$ αρχίζουμε και «τρέχουμε» προς κάθε δυνατή κατεύθυνση και σταματάμε μόνο αν «πέσουμε» πάνω σε κάποιον τοίχο του δωματίου μας (δηλαδή, κάποια πλευρά του $P$). Επομένως, το $\phi_P(A)$ θα μοιράζεται κάποιες πλευρές (ή κάποια τμήματά τους) με το $P$ ενώ, σε περιπτώσεις σαν κι αυτές που φαίνονται στο Σχήμα \ref{fig:6}, με τις διακεκομμένες γραμμές, τον ρόλο των πλευρών του $\phi_P(A)$ τον παίζουν οι ίδιες οι διακεκομμένες γραμμές.
	
	Κάτι που έχει ενδιαφέρον να εξετάσουμε τώρα και που θα μας βοηθήσει να καταλάβουμε πώς μοιάζει το $\phi_P(A)$ είναι το ποιες \textit{κορυφές} του $P$ φωτίζονται από το $A$. Ας πάρουμε μία παραλλαγή του δωματίου που εξετάσαμε προηγουμένως, όπως φαίνεται στο Σχήμα \ref{fig:8}. Εκεί έχουμε στραβώσει κάπως το ένα «πόδι» του «Π» που σχημάτιζε το δωμάτιο, όπως επίσης έχουμε στραβώσει λίγο και κάποιες άλλες πλευρές του δωματίου (γιατί το να έχουμε ένα δωμάτιο μόνο με ορθές γωνίες ήταν μία «ριψοκίνδυνη» επιλογή για τη γενικότητα των συλλογισμών μας).
	
	\begin{figure}[!htb]
		\centering
		\begin{tikzpicture}
		\fill[red!10] (0,13/6) -- (1,2.5) -- (3,1.5) -- (3.5,0) -- (4,0) -- (4,4) -- (0,3) -- cycle;
		\draw[thick] (0,0) -- (0,3) -- (4,4) -- (4,0) -- (3,0) -- (3,1.5) -- (1,2.5) -- (1,0) -- cycle;
		\node[circle, inner sep=1.5pt, draw=black, fill=black, label={right}:{$A$}](a) at (2.5,3){};
		\draw[dashed] (a) -- (3.5,0);
		\draw[dashed] (a) -- (0,13/6);
		\node[circle, draw=red, fill=red, inner sep=1.5pt] at (4,0){};
		\node[circle, draw=red, fill=red, inner sep=1.5pt] at (4,4){};
		\node[circle, draw=red, fill=red, inner sep=1.5pt] at (0,3){};
		\node[circle, draw=red, fill=red, inner sep=1.5pt] at (3,1.5){};
		\node[circle, draw=red, fill=red, inner sep=1.5pt] at (1,2.5){};
		\node[circle, draw=red, fill=white, inner sep=1.5pt] at (1,0){};
		\node[circle, draw=red, fill=white, inner sep=1.5pt] at (0,0){};
		\node[circle, draw=red, fill=white, inner sep=1.5pt] at (3,0){};
		\end{tikzpicture}
		\caption{'Οταν τα πράγματα στραβώνουν\ldots}
		\label{fig:8}
	\end{figure}

	Στο Σχήμα \ref{fig:8} έχουμε σημειώσει με κόκκινα «γεμιστά» κυκλάκια εκείνες τις κορυφές του $P$ τις οποίες η λάμπα μας φωτίζει και με κόκκινα «άδεια» κυκλάκια εκείνες τις οποίες δε φωτίζει. Φαίνεται, από το εν λόγω σχήμα, να μπορούμε να διατυπώσουμε δύο εύλογες εικασίες για τη φύση του $\phi_P(A)$:
	\begin{enumerate}
		\item Αν $x,y$ είναι δύο διαδοχικές κορυφές του $P$ με $x,y\in\phi_P(A)$ τότε και η πλευρά $[x,y]\subseteq\phi_P(A)$.
		\item Αν $x,y$ είναι δύο διαδοχικές κορυφές του $P$ με $x\in\phi_P(A)$ και $y\not\in\phi_P(A)$ τότε $[x,y]\not\subseteq\phi_P(A)$ και υπάρχει $z\in[x,y]$ έτσι ώστε\footnote{Δανειζόμαστε τον συμβολισμό των διαστημάτων πραγματικών αριθμών για να συμβολίσουμε τα ευθύγραμμα τμήματα του $\R^2$. 'Ετσι, όπως είπαμε, το $[x,y]$ συμβολίζει ένα ευθύγραμμο τμήμα με άκρα $x,y$ συμπεριλαμβάνοντας και τα άκρα του, ενώ το $(x,y)$ συμβολίζει το ίδιο ευθύγραμμο τμήμα, απλά χωρίς τα άκρα του. Αναλόγως και για τα $(x,y]$ και $[x,y)$.} $[x,z]\subseteq\phi_P(A)$ και $(z,y]\cap\phi_P(A)=\varnothing$.
	\end{enumerate}

	Οι δύο παραπάνω εικασίες, αν αληθεύουν, μας λένε πολλά για τη δομή του $\phi_P(A)$ καθώς μπορούν να μας βοηθήσουν να καθορίσουμε αλγοριθμικά και σε πεπερασμένο χρόνο ποιες είναι οι κορυφές και οι πλευρές του. Επομένως, θα τις σκαλίσουμε λίγο ακόμα. Η πρώτη εικασία μας λέει κάτι απλό: αν η λάμπα μας μπορεί να φωτίσει δύο διαδοχικές «γωνίες» του δωματίου μας τότε θα μπορεί να φωτίσει και τον τοίχο που τις ενώνει. Αυτό, από την εμπειρία μας φαντάζει προφανές, οπότε ίσως έχουμε την τύχη με το μέρος μας. Ας πάρουμε, λοιπόν, δύο διαδοχικές κορυφές $x,y\in P$ κι ας υποθέσουμε ότι αμφότερες ανήκουν και στο $\phi_P(A)$, δηλαδή ότι $[A,x]\subseteq P$ και ότι $[A,y]\subseteq P$. Το ζητούμενο είναι να αποδείξουμε ότι $[x,y]\subseteq\phi_P(A)$, δηλαδή ότι $z\in[x,y]\Rightarrow z\in\phi_P(A)$ ή, χρησιμοποιώντας την Εξίσωση \eqref{eq:light_pa}, ότι:
	\begin{equation}
	[A,z]\subseteq P.
	\end{equation}
	Θα προχωρήσουμε με απαγωγή σε άτοπο. 'Εστω, λοιπόν, ότι υπάρχει κάποιο $z\in[x,y]$ τέτοιο ώστε $z\not\in\phi_P(A)$, δηλαδή τέτοιο ώστε $[A,z]\not\subseteq P$. Αυτό, με τη σειρά του, σημαίνει ότι υπάρχει κάποιο $u\in[A,z]$ τέτοιο ώστε $u\not\in P$ και, μάλιστα, επειδή $A,z\in P$, έπεται ότι $u\in(A,z)$. Μιας και μαζεύτηκαν κάπως πολλά, ας δούμε λίγο το Σχήμα \ref{fig:9} για να ξεκαθαρίσουμε την κατάσταση.
	
	\begin{figure}[!htb]
		\centering
		\begin{tikzpicture}
		\node[circle, draw=red, fill=red, inner sep=1.5pt, label={below}:{$A$}](a) at (0,-1){};
		\node[circle, draw=red, fill=red, inner sep=1.5pt, label={right}:{$x$}](x) at (2,3){};
		\node[circle, draw=red, fill=red, inner sep=1.5pt, label={left}:{$y$}](y) at (-1,1){};
		\node[circle, draw=red, fill=red, inner sep=1.5pt, label={above}:{$z$}](z) at (.5,2){};
		\node[circle, draw=red, fill=white, inner sep=1.5pt, label={left}:{$u$}](u) at (1/6,0){};
		\draw[thick, red!30] (y) -- (z);
		\draw[thick, red!30] (x) -- (z);
		\draw[thick, dashed, red!30] (a) -- (x);
		\draw[thick, dashed, red!30] (a) -- (y);
		\draw[dashed] (a) -- (u);
		\draw[dashed] (u) -- (z);
		\end{tikzpicture}
		\caption{Τα σημεία που μας ενδιαφέρουν.}
		\label{fig:9}
	\end{figure}

	Τα «γεμιστά» σημεία του σχήματος ανήκουν όλα τους στο $P$ ενώ το $u$, όπως είπαμε, δεν ανήκει. Επίσης, όλα τα κόκκινα ευθύγραμμα τμήματα (είτε διακεκομμένα είτε όχι) περιέχονται επίσης στο $P$. Τώρα, εμείς θέλουμε κάπως να καταλήξουμε σε ένα άτοπο. Ως τώρα έχουμε αξιοποιήσει τις περισσότερες υποθέσεις μας, πέρα από το γεγονός ότι οι $x,y$ είναι \textit{διαδοχικές} κορυφές του $P$, το οποίο σημαίνει κι ότι κάθε σημείο του $[x,y]$ ανήκει στο $P$ (αφού το $[x,y]$ είναι πλευρά του $P$). Τώρα, αφού τα $A,z\in P$ αλλά το $u\not\in P$ πρέπει με κάποιον τρόπο να βρίσκεται το $u$ «περικυκλωμένο» από πλευρές του $P$. Για την ακρίβεια θα πρέπει το σύνορο του $P$ (η τεθλασμένη γραμμή που το περιγράφει, δηλαδή), να τέμνει το $[A,z]$ σε τουλάχιστον άλλα δύο σημεία πέρα από το $z$, έστω $v,w$, έτσι ώστε το $(v,w)$ να βρίσκεται στο εξωτερικό του $P$, δηλαδή $(v,w)\cap P=\varnothing$ (βλ. και Σχήμα \ref{fig:10}).
	
	\begin{figure}[!htb]
		\centering
		\begin{tikzpicture}
		\node[circle, draw=red, fill=red, inner sep=1.5pt, label={below}:{$A$}](a) at (0,-1){};
		\node[circle, draw=red, fill=red, inner sep=1.5pt, label={right}:{$x$}](x) at (2,3){};
		\node[circle, draw=red, fill=red, inner sep=1.5pt, label={left}:{$y$}](y) at (-1,1){};
		\node[circle, draw=red, fill=red, inner sep=1.5pt, label={above}:{$z$}](z) at (.5,2){};
		\node[circle, draw=red, fill=red, inner sep=1.5pt, label={right}:{$v$}](v) at (1/12,-.5){};
		\node[circle, draw=red, fill=red, inner sep=1.5pt, label={left}:{$w$}](w) at (7/24,.75){};
		\node[circle, draw=red, fill=white, inner sep=1.5pt, label={left}:{$u$}](u) at (1/6,0){};
		\draw[thick, red!30] (y) -- (z);
		\draw[thick, red!30] (x) -- (z);
		\draw[thick, dashed, red!30] (a) -- (x);
		\draw[thick, dashed, red!30] (a) -- (y);
		\draw[dashed] (a) -- (v);
		\draw[blue!50] (v) -- (u);
		\draw[blue!50] (u) -- (w);
		\draw[dashed] (w) -- (z);
		\end{tikzpicture}
		\caption{Μία «τρύπα» στο $P$.}
		\label{fig:10}
	\end{figure}

	Ωστόσο, εδώ θα πρέπει να αρχίσουμε κάπως να προβληματιζόμαστε. 'Ολα τα «γεμιστά» σημεία πλην του $A$ βρίσκονται στο σύνορο του πολυγώνου $P$ και, μάλιστα, τα $x,y$ είναι και κορυφές του. Ας πούμε τώρα ότι στεκόμαστε στο σημείο $x$ και θέλουμε να περπατήσουμε κατά μήκος του συνόρου του $P$. 'Εχουμε δύο επιλογές: η μία είναι να κινηθούμε προς το $y$ και η άλλη μακριά από αυτό. Ας επιλέξουμε τη δεύτερη. Τότε, επειδή το $[x,y]$ είναι η τελευταία πλευρά του $P$ που θα περπατήσουμε κι επειδή τα $v,w$ δε βρίσκονται στο $[x,y]$ θα πρέπει να περάσουμε από αυτά πριν φτάσουμε στο $y$. Για παράδειγμα η διαδρομή μας μπορεί να μοιάζει κάπως σαν κι αυτή που φαίνεται στο Σχήμα \ref{fig:11}, δεδομένου ότι ό,τι εμφανίζεται ως κόκκινο (είτε είναι σημείο είτε γραμμή) πρέπει να βρίσκεται μέσα στο πολύγωνό μας, εκτός από το σημείο $u$ και το μπλε ανοικτό ευθύγραμμο τμήμα $v,w$.
	
	\begin{figure}[!htb]
		\centering
		\begin{tikzpicture}
		\node[circle, draw=red, fill=red, inner sep=1.5pt, label={below}:{$A$}](a) at (0,-1){};
		\node[circle, draw=red, fill=red, inner sep=1.5pt, label={right}:{$x$}](x) at (2,3){};
		\node[circle, draw=red, fill=red, inner sep=1.5pt, label={left}:{$y$}](y) at (-1,1){};
		\node[circle, draw=red, fill=red, inner sep=1.5pt, label={above}:{$z$}](z) at (.5,2){};
		\node[circle, draw=red, fill=red, inner sep=1.5pt, label={right}:{$v$}](v) at (1/12,-.5){};
		\node[circle, draw=red, fill=red, inner sep=1.5pt, label={left}:{$w$}](w) at (7/24,.75){};
		\node[circle, draw=red, fill=white, inner sep=1.5pt, label={left}:{$u$}](u) at (1/6,0){};
		\draw[thick, red!30] (y) -- (z);
		\draw[thick, red!30] (x) -- (z);
		\draw[thick, dashed, red!30] (a) -- (x);
		\draw[thick, dashed, red!30] (a) -- (y);
		\draw[dashed] (a) -- (v);
		\draw[blue!50] (v) -- (u);
		\draw[blue!50] (u) -- (w);
		\draw[dashed] (w) -- (z);
		\begin{pgfonlayer}{bg}
		\begin{scope}[decoration={footprints, foot length=4pt, stride length=10pt, foot of=human}]
			\draw[decorate, gray!50, fill=gray!10] (x) to[bend left] (2.4,2) to[bend left] (w) to[bend right]  (-.75,0) to[bend right] (v) to[bend left] (1,-1) to[bend left] (0,-2) to[bend left] (-1.6,0) to[bend left] (y);
		\end{scope}
		\end{pgfonlayer}
		\end{tikzpicture}
		\caption{Περπατώ, περπατώ εις το δάσος, όταν ο Gauss δεν είν' εδώ\ldots}
		\label{fig:11}
	\end{figure}

	Αυτός μας ο περίπατος έχει ένα χαρακτηριστικό: σε κάποια στιγμή θα χρειαστεί να περάσουμε πάνω από τουλάχιστον ένα από τα ευθύγραμμα τμήματα $[A,x]$ ή $[A,y]$ για να «μπούμε» στο τρίγωνο $\triangle Axy$ και άλλη μία φορά για να «βγούμε» από αυτό. Αυτό είναι απαραίτητο καθώς δε μπορούμε, ξεκινώντας από το $x$, να «μπούμε» κατευθείαν στο τρίγωνο $\triangle Axy$ καθώς όλα τα σημεία του $[A,x]$ πρέπει να ανήκουν στο $P$. 'Ετσι, πρέπει να ξεκινήσουμε είτε κινούμενοι επί της $[A,x]$ είτε έξω από αυτήν, όπως στο Σχήμα \ref{fig:11}. Σε κάποια περίπτωση, όπως κι αν έχει η κατάσταση, θα χρειαστεί να «μπούμε» στο τρίγωνο για να περάσουμε από τα σημεία $u$ και $w$, οπότε θα κόψουμε μία από τις πλευρές $[A,x]$ ή $[A,y]$. 'Επειτα, για να καταφέρουμε να συμπεριλάβουμε το σημείο $A$ στο $P$ θα πρέπει είτε να βγούμε από το τρίγωνο είτε, έστω, να κινηθούμε επί των πλευρών του και να περάσουμε τόσο από το $A$ όσο και από το $y$. Ωστόσο, το σημείο στο οποίο θα ξανασυναντήσουμε το τις πλευρές του τριγώνου δε θα είναι αυτό από το οποίο «μπήκαμε»\footnote{Διότι δεν μπορεί το σύνορο ενός πολυγώνου να τέμνει τον εαυτό του.}, συνεπώς, κάποια σημεία των πλευρών θα τα αφήσουμε αναγκαστικά εκτός του $P$, όπως φαίνεται και στο Σχήμα \ref{fig:11}.
	
	Σαφώς, τα παραπάνω δεν είναι όλα τους διατυπωμένα με αυστηρά μαθηματικά, αλλά είναι αρκετά εύκολο (πλην όμως, κουραστικό) να τα αποτυπώσει κανείς αυστηρά. Η διαίσθηση ωστόσο παραμένει η ίδια και το συμπέρασμα είναι ότι κάποιο σημείο επί των $[A,x]$ ή $[A,y]$ πρέπει να μείνει εκτός του $P$, άτοπο! Συνεπώς, δεν υπάρχει σημείο της $[x,y]$ που να μην ανήκει στην $\phi_P(A)$ και άρα $[x,y]\subseteq\phi_P(A)$ όπως είχαμε εικάσει.
	
	Πριν περάσουμε στη δεύτερη εικασία, αξίζει να σταθούμε λίγο ακόμα σε αυτήν εδώ που μόλις αποδείξαμε. Αρχικά, θα τη διατυπώσουμε υπό μορφή πρότασης για να μπορούμε να την καμαρώνουμε:
	\begin{prop}[Εικασία 1\textsuperscript{η}]\label{prop:conjecture_1}
		Αν $P$ είναι ένα πολύγωνο, $A\in P$ και $\phi_P(A)$ είναι η «φωτιζόμενη» περιοχή του πολυγώνου από το $A$ όπως ορίζεται από την Εξίσωση \eqref{eq:light_pa} και αν $x,y\in P$ είναι δύο διαδοχικές κορυφές του $P$ έτσι ώστε $x,y\in\phi_P(A)$ τότε ισχύει και ότι $[x,y]\subseteq\phi_P(A)$.
	\end{prop}

	\begin{figure}[!htb]
		\centering
		\begin{tikzpicture}
		\fill[red!10] (0,13/6) -- (1,2.5) -- (3,1.5) -- (3.5,0) -- (4,0) -- (4,4) -- (0,3) -- cycle;
		\fill[red!20] (0,3) -- (4,4) -- (2.5,3) -- cycle;
		\fill[red!20] (4,0) -- (4,4) -- (2.5,3) -- cycle;
		\fill[red!20] (3,1.5) -- (1,2.5) -- (2.5,3) -- cycle;
		\draw[thick] (0,0) -- (0,3) -- (4,4) -- (4,0) -- (3,0) -- (3,1.5) -- (1,2.5) -- (1,0) -- cycle;
		\node[circle, inner sep=1.5pt, draw=black, fill=black, label={right}:{$A$}](a) at (2.5,3){};
		\node[circle, draw=red, fill=red, inner sep=1.5pt] (x) at (4,0){};
		\node[circle, draw=red, fill=red, inner sep=1.5pt] (y) at (4,4){};
		\node[circle, draw=red, fill=red, inner sep=1.5pt] (z) at (0,3){};
		\node[circle, draw=red, fill=red, inner sep=1.5pt] (u) at (3,1.5){};
		\node[circle, draw=red, fill=red, inner sep=1.5pt] (v) at (1,2.5){};
		\node[circle, draw=red, fill=white, inner sep=1.5pt] (w) at (1,0){};
		\node[circle, draw=red, fill=white, inner sep=1.5pt] (s) at (0,0){};
		\node[circle, draw=red, fill=white, inner sep=1.5pt] (t) at (3,0){};
		\draw[red] (a) -- (u);
		\draw[red] (a) -- (v);
		\draw[red] (a) -- (x);
		\draw[red] (a) -- (y);
		\draw[red] (a) -- (z);
		\end{tikzpicture}
		\caption{'Οταν τα πράγματα στραβώνουν\ldots}
		\label{fig:12}
	\end{figure}
	
	Ωραία, τώρα που βγάλαμε την τυπική αυτή υποχρέωση, ας δούμε κάτι πιο ζουμερό. Ας παρατηρήσουμε λίγο το Σχήμα \ref{fig:12} στο οποίο έχουμε σχεδιάσει στο παραμορφωμένο αγαπημένο μας δωμάτιο κάποιες «φωτεινές ακτίνες» που ενώνουν τη λάμπα μας με τις γωνίες του δωματίου που αυτή μπορεί να φωτίσει άμεσα. 'Οπως μπορεί κανείς εύκολα να παρατηρήσει, από το συμπέρασμα της Πρότασης \ref{prop:conjecture_1} και το γεγονός ότι πέρα από την πλευρά $[x,y]$ (με τους όρους της Πρότασης \ref{prop:conjecture_1}), και τα ευθύγραμμα τμήματα $[A,x]$ και $[A,y]$ περιέχονται στο $\phi_P(A)$, έπεται ότι όλο το τρίγωνο $\triangle Axy$ περιέχεται στο $\phi_P(A)$. Αυτό είναι εύκολο να το αποδείξει κανείς, αφού, αν $a\in\phi_P(A)$ και $b\in[A,a]$ τότε είναι άμεσο από την Εξίσωση \eqref{eq:light_pa} ότι $b\in\phi_P(A)$, αφού $[A,b]\subseteq[A,a]$ (επειδή $b\in[A,a]$) και $[A,a]\subseteq P$ (αφού $a\in\phi_P(A)$). Συνεπώς, $[A,b]\subseteq P$ και άρα $b\in\phi_P(A)$.
	
	Από την παραπάνω ιδιότητα όμως παίρνουμε άμεσα αυτό που θέλουμε, καθώς αν $a\in\triangle Axy$ τότε προεκτείνουμε το $[A,a]$ μέχρι να συναντήσουμε στο σημείο $a'$ την $[x,y]$. Τώρα, επειδή το $a'$ θα ανήκει στην $[x,y]$, θα ανήκει και στην $\phi_P(A)$. Από την άλλη, εκ κατασκευής $a\in[A,a']$ οπότε, από τα παραπάνω έπεται ότι $a\in\phi_P(A)$ κι επειδή το $a\in\triangle Axy$ ήταν ένα αυθαίρετο σημείο του τριγώνου, αυτό ισχύει για κάθε σημείο του, άρα $\triangle Axy\subseteq\phi_P(A)$. Σχηματικά, όλο αυτό το απλό επιχείρημα αποτυπώνεται στο Σχήμα \ref{fig:13}.
	
	\begin{figure}[!htb]
		\centering
		\begin{tikzpicture}
		\node[circle, draw=red, fill=red, inner sep=1.5pt, label={below}:{$A$}](A) at (0,-1){};
		\node[circle, draw=red, fill=red, inner sep=1.5pt, label={right}:{$x$}](x) at (2,3){};
		\node[circle, draw=red, fill=red, inner sep=1.5pt, label={left}:{$y$}](y) at (-1,1){};
		\node[circle, draw=red, fill=red, inner sep=1.5pt, label={above}:{$a'$}](a) at (.5,2){};
		\node[circle, draw=red, fill=red, inner sep=1.5pt, label={left}:{$a$}](ap) at (1/6,0){};
		\draw[thick, red!30] (y) -- (a);
		\draw[thick, red!30] (x) -- (a);
		\draw[thick, dashed, red!30] (A) -- (x);
		\draw[thick, dashed, red!30] (A) -- (y);
		\draw[dashed] (A) -- (ap);
		\draw[dashed] (ap) -- (a);
		\end{tikzpicture}
		\caption{Η απόδειξη μίας απλής ιδιότητας.}
		\label{fig:13}
	\end{figure}

	Αυτή η απλή ιδιότητα μπορεί να γραφτεί σαν ένα κομψό πόρισμα της Πρότασης \ref{prop:conjecture_1} ως εξής:
	\begin{cor}[της Πρότασης \ref{prop:conjecture_1}]\label{cor:conjecture_1}
		Αν $P$ είναι ένα πολύγωνο, $A\in P$ και $\phi_P(A)$ είναι η «φωτιζόμενη» περιοχή του πολυγώνου από το $A$ όπως ορίζεται από την Εξίσωση \eqref{eq:light_pa} τότε αν $x,y\in P$ είναι δύο διαδοχικές κορυφές του $P$ έτσι ώστε $x,y\in\phi_P(A)$ τότε ισχύει και ότι $\triangle Axy\subseteq\phi_P(A)$.
	\end{cor}
	
	Το παραπάνω πόρισμα, αν το χρησιμοποιήσουμε στο παραμορφωμένο αγαπημένο μας δωμάτιο, μας εγγυάται ότι τα εντόνως κόκκινα τρίγωνα που φαίνονται στη Σχήμα \ref{fig:12} φωτίζονται πράγματι από τη λάμπα μας που έχει τοποθετηθεί στο $A$. 'Ηδη, λοιπόν, αν και δεν έχουμε καταφέρει, όπως βλέπετε, να περιγράψουμε όλο το $\phi_P(A)$ ακόμα, έχουμε καταφέρει να καλύψουμε ένα αρκετά μεγάλο μέρος του δρόμου μας προς αυτόν τον στόχο. Για την ακρίβεια, έχουμε αποδείξει ότι τα τριγωνάκια που καθορίζονται από τη θέση της λάμπας μας και τις γωνίες του δωματίου που μπορεί αυτή να φωτίσει φωτίζονται κι αυτά εξ ολοκλήρου\footnote{Ναι, λες και περιμέναμε τα μαθηματικά να μας το πουν, δεν το βλέπαμε και με τα μάτια μας αυτό\ldots}. Τώρα μένει να εξερευνήσουμε αυτές τις αχνές περιοχές που φαίνονται στο Σχήμα \ref{fig:12}.
	
	\section{Μήπως ξεχάσαμε κάτι;}
	Ναι, ξεχάσαμε να διερευνήσουμε τη δεύτερη εικασία που διατυπώσαμε παραπάνω. Βασικά, δεν το ξεχάσαμε, απλώς η προηγούμενη ενότητα έκλεισε νοηματικά, οπότε θα συνεχίσουμε τη διερεύνησή μας σ' αυτήν εδώ την ενότητα.
\end{document}